\chapter{PACT-PAIR: Dyadic Evaluation}
\label{chapter:benchmark_design}

This chapter introduces the PACT-PAIR benchmark (Privacy in Agent Cross-boundary Transactions --- Pairwise Assessment of Information Risk), a 600-task evaluation suite probing one requesting agent against one target agent across a single privacy boundary.

\begin{tcolorbox}[colback=blue!3!white, colframe=blue!40!black, title={\textbf{RQ1: What moves the frontier?}}]
We vary policy specificity from none to category enumeration.
\end{tcolorbox}
\begin{tcolorbox}[colback=blue!3!white, colframe=blue!40!black, title={\textbf{RQ2: Is the frontier stable over time?}}]
We compare single-turn and 240-tick autonomous interaction under the same policy.
\end{tcolorbox}
\begin{tcolorbox}[colback=blue!3!white, colframe=blue!40!black, title={\textbf{RQ3: Is the frontier the same for everyone?}}]
We fix D2 and vary only the requester's relationship.
\end{tcolorbox}


% ============================================================================
\section{Setup}
\label{sec:pact_pair_setup}
% ============================================================================

\textbf{Persona.} The target agent (Alex) is a co-founder/CTO with hundreds of files and structured state spanning six sensitivity categories. The requesting agent (Tina) communicates through the cross-boundary protocol.

\textbf{Task composition.} 600 tasks across three tracks: Files QA (200), States QA (200), and Actions (200). Each track is balanced: 100 utility tasks and 100 security tasks. Each task is paired with gold key facts for automated evaluation.

\textbf{Models.} Six models from three providers: gpt-5-mini (ablation backbone), gpt-5.5, gpt-5.4-mini, gpt-5.4, Kimi K2, DeepSeek V3.

\textbf{Evaluation.} Two-pass pipeline: LLM judge + string match with OR-adjudication. Action verdicts use database-snapshot diffs. Manual audit ($\kappa = 0.96$, $n=60$) confirmed zero false positives.


% ============================================================================
\section{RQ1: What Moves the Frontier}
\label{sec:rq1_findings}
% ============================================================================

\begin{table}[h]
  \caption{Single-step files QA headline metrics (gpt-5-mini, 2 replications, OR-pooled).}
  \label{tab:rq1_files_headline}
  \centering
  \small
  \begin{tabular}{lcccc}
    \toprule
    \textbf{Defence} & \textbf{Utility} & \textbf{Leak Rate} & \textbf{Refuse Rate} & \textbf{$n$ (security)} \\
    \midrule
    D0 (no policy) & 78.0\% & 83.0\% & 0.0\% & 200 \\
    D1 (generic) & 78.5\% & 81.5\% & 3.0\% & 200 \\
    D2 (category deny list) & 77.0\% & 14.0\% & 71.5\% & 200 \\
    \bottomrule
  \end{tabular}
\end{table}

\textbf{The specificity threshold.} D1 is statistically indistinguishable from D0 (McNemar $p = \text{n.s.}$). D0 versus D2: discordant ratio 70:1 ($p < 0.001$). The D0$\to$D2 gap is 69 percentage points of leak reduction with only 1pp utility cost. Below a specificity threshold, governance has no measurable effect. Category enumeration converts privacy decisions from open-ended judgment into pattern matching---a task models perform reliably.

\textbf{Per-category hierarchy.} Sensitive work leaks most (28.3\% under D2) because organisational data falls between individual-level deny categories. Personal finance (4.0\%) and health (12.5\%) are near-zero because they are unambiguously individual.

\textbf{Surface asymmetry.} D2 costs only 1pp utility on files but 37pp on structured state. Terse state queries share vocabulary with protected categories, triggering 26.5\% false-refusal.

\textbf{Cross-model generality.} The specificity threshold holds across all six models (D2 reduces leakage by 69--91pp in every model).

\textbf{Action safety.} D2 is Pareto-superior: 93.5\% safety with 61.0\% utility. D1 is Pareto-\textit{inferior}: identical safety to D0 while reducing utility by 17.5pp. D2 achieves 100\% block rate on destructive actions.

\begin{quote}
\textbf{Takeaway (RQ1):} The frontier only moves when you name the boundary. Category enumeration crosses a specificity threshold below which governance has no measurable effect. This threshold is universal across six models from three providers.
\end{quote}


% ============================================================================
\section{RQ2: Is the Frontier Stable Over Time}
\label{sec:rq2_findings}
% ============================================================================

\begin{table}[h]
  \caption{Multi-turn D2 security (gpt-5-mini, 10 splits $\times$ 240 ticks, $n=191$).}
  \label{tab:rq2_category}
  \centering
  \small
  \begin{tabular}{lcccc}
    \toprule
    \textbf{Category} & \textbf{Refuse} & \textbf{Msg.~Leak} & \textbf{Failed Att.} & \textbf{Glob.~Leak} \\
    \midrule
    Sensitive work         & 36.8 & 22.8 & 40.4 & 51.7 \\
    Personal finance       & 75.0 &  4.5 & 20.5 & 34.7 \\
    Personal health        & 89.7 &  2.6 &  7.7 & 20.0 \\
    Personal relationships & 66.7 & 15.7 & 17.6 & 39.2 \\
    \midrule
    \textbf{All sensitive} & \textbf{64.4} & \textbf{12.6} & \textbf{23.0} & \textbf{38.0} \\
    \bottomrule
  \end{tabular}
\end{table}

\textbf{Bounded erosion, not collapse.} D2's multi-turn message leak rate (12.6\%) is comparable to single-turn (14.0\%). But the global leak rate (38.0\%) is 3$\times$ the message leak because sensitive facts surface incidentally in responses to unrelated questions.

\textbf{Adaptive strategies.} Strategy \textit{combinations} are more dangerous than individual strategies. The wedding cascade case study required 12 failed attempts before a three-strategy combination (business justification + scope expansion + constrained format) succeeded. After the breach, the policy recovered on six subsequent probes.

\textbf{D3/D4/D5 marginal improvement.} All three defences provide only 4--11pp leak reduction at 6--13pp utility cost atop D2. Even the best (D5) still leaks 27.5\% over 240 ticks.

\begin{quote}
\textbf{Takeaway (RQ2):} Multi-turn erodes the frontier through incidental channels the policy never evaluates. The global leak rate is 3$\times$ the per-message rate. Single-turn benchmarks systematically undercount real-world privacy risk.
\end{quote}


% ============================================================================
\section{RQ3: Is the Frontier the Same for Everyone}
\label{sec:rq3_findings}
% ============================================================================

\begin{table}[h]
  \caption{Relationship-conditioned D2 results (gpt-5-mini, 400 QA + 200 actions per requester).}
  \label{tab:rq3_relationship}
  \centering
  \small
  \begin{tabular}{llccc}
    \toprule
    \textbf{Requester} & \textbf{Role} & \textbf{Leak}$\downarrow$ & \textbf{Over-Refusal}$\downarrow$ & \textbf{Utility}$\uparrow$ \\
    \midrule
    Tina   & Colleague       & 1.7\%  & 40\%  & 57\% \\
    Marcus & CEO delegate    & 3.3\%  & 59\%  & 49\% \\
    Jordan & Close friend    & 9.2\%  & 86\%  & 58\% \\
    Dana   & Investor/board  & 7.5\%  & 31\%  & 70\% \\
    \bottomrule
  \end{tabular}
\end{table}

\textbf{Only sensitive work is relationship-dependent.} Personal finance, health, and relationships form a hard floor regardless of who asks. Sensitive work shows a clear gradient: Tina 0\%, Marcus 20\%, Dana 26\%.

\textbf{Over-refusal is the dominant failure.} Jordan's agent is refused on 86\% of items it \textit{should} access. D2's category-level rules take precedence over relationship inference: a close friend asking about emergency medication is refused identically to a stranger.

\textbf{Read trust $\neq$ write trust.} Dana (investor) has highest QA leakage but also highest action safety. Jordan (friend) has most over-refusal yet lowest action safety. Trust decomposes into orthogonal read/write axes.

\begin{quote}
\textbf{Takeaway (RQ3):} A single policy cannot serve all requesters optimally. Over-refusal, not leakage, is the dominant practical failure. These findings motivate per-requester retrieval scope (MCC).
\end{quote}


% ============================================================================
\section{Failure Taxonomy}
\label{sec:failure_taxonomy}
% ============================================================================

D2 residual leaks fall into four mechanisms: \textbf{category boundary ambiguity} (61\%)---organisational data not covered by individual-level deny list; \textbf{leaked-outside-message} (16\%)---refusal text reveals protected facts; \textbf{company-benefit confusion} (7\%)---benefits treated as company info; \textbf{stochastic failures} (16\%)---identical inputs, different outputs across replications. The dominant failure (61\%) cannot be fixed by longer policies. It requires moving enforcement from reasoning to infrastructure: controlling what data the agent can access rather than trusting it to decide what to disclose.
